% Physical unit 023; accepted Session 02.
\GordanUnitStart{023}{printed page 18}
\GordanTopFolio{18}
\begingroup\scriptsize
\[
\begin{aligned}
={}&(11)_2\omega^2\left\{\left({dl\over d\psi_1}\right)_2\left({dl\over dl_0}\right)-
\left({dl\over d\psi_1}\right)\left({dl\over dl_0}\right)_2\right\}
\left\{\left({dl\over d\psi_1}\right)_2\left({dl\over dl_0}\right)_1-
\left({dl\over d\psi_1}\right)_1\left({dl\over dl_0}\right)_2\right\}\\
&\times\left\{\left({dl\over d\psi_1}\right)_1\left({dl\over dl_0}\right)-
\left({dl\over d\psi_1}\right)\left({dl\over dl_0}\right)_1\right\}
\left\{\left({dl\over d\psi_1}\right)_2\left({dl''\over dl_0}\right)_2-
\left({dl''\over d\psi_1}\right)_2\left({dl\over dl_0}\right)_2\right\}
\end{aligned}
\]
\endgroup
et posito (cap. IV. \S\ 7):
\[
 {dl\over d\psi_1}={\sin l_0\sin\varphi\over v};
 \qquad {dl\over dl_0}={u\over v}
\]
\begingroup\scriptsize
\[
\begin{aligned}
U={}&(11)_2\omega^2
\left({\sin l_0\sin\varphi_2\over v_2}{u\over v}-{\sin l_0\sin\varphi\over v}{u_2\over v_2}\right)
\left({\sin l_0\sin\varphi_2\over v_2}{u_1\over v_1}-{\sin l_0\sin\varphi_1\over v_1}{u_2\over v_2}\right)\\
&\times\left({\sin l_0\sin\varphi_1\over v_1}{u\over v}-{\sin l_0\sin\varphi\over v}{u_1\over v_1}\right)
\left\{\left({\sin l_0\sin\varphi_2\over v_2}\right)\left({u\over v}\right)'_2-
\left({\sin l_0\sin\varphi\over v}\right)'_2\left({u_2\over v_2}\right)\right\}.
\end{aligned}
\]
\endgroup
Quodsi reputamus, posito $\psi=\psi_1$, $u$ evanescere, id est
\[
 u_1=0
\]
esse; sequitur:
\begingroup\scriptsize
\[
\begin{aligned}
U={}&(11)_2\omega^2\sin^4l_0
\left\{{u\sin\varphi_2\over vv_2}-{u_2\sin\varphi\over vv_2}\right\}
\left\{-{u_2\sin\varphi_1\over v_1v_2}\right\}
\left\{{u\sin\varphi_1\over vv_1}\right\}\\
&\times\left\{{u_2\over v_2}\left({\sin\varphi\over v}\right)'_2-{\sin\varphi_2\over v_2}\left({u\over v}\right)'_2\right\}\\
={}&(11)_2\omega^2{uu_2\sin^4l_0\sin\varphi_1\sin\varphi_2\over v^2v_1^2v_2^2}
\{u\sin\varphi_2-u_2\sin\varphi\}\,{\sin\varphi_2^2\over v_2^2}\cdot\left\{{u\over \sin\varphi}\right\}'_2\\
={}&(11)_2\omega^2{uu_2\sin^4l_0\sin\varphi_1^2\sin\varphi_2^2\over v^2v_1^2v_2^4}
\{u\sin\varphi_2-u_2\sin\varphi\}\left\{-{\cos^2l\over \cos l_0\sin^2\varphi_2}\right\},\\
U={}&(11)_2\sin^4l_0{\sin^2\varphi_1\sin\varphi_2^2\sin^2\varphi\cos^2l\over \cos l_0}{1\over v^2}{u_2\over v_1^2}{1\over v_2^4}
\left(-{u_2\over \sin\varphi_2}\right)\cdot{u\over \sin\varphi}\left\{{u\over \sin\varphi}-{u_2\over \sin\varphi_2}\right\}.
\end{aligned}
\]
\endgroup

\begin{center}\S\ 14.\end{center}

Hactenus cum $\psi_2$, inter $\psi$ et $\psi_1$ situm, non accuratius tamen definitum, consideravimus, $\psi_2$ jam angustioribus coërcendum esse videatur limitibus.

Quoniam enim
\[
 {u\over \sin l_0\sin\varphi},
\]
ut supra demonstravimus, pro ratione crescentis longitudinis deinceps numeros a $0$ ad $-\infty$ induit; $\psi_2$ tam prope $\psi_1$ situm sumi potest, ut $\psi$ ad valorem $\psi_2$ producatur, priusquam
\[
 {u\over \sin\varphi}
\]
ad $-\infty$ adducatur.

Quare quoniam
\[
 {-u_2\over \sin\varphi_2}>0
\]
